\lysh{35143}{2}

\begin{alist}
\item Από τα δεδομένα της άσκησης βρίσκουμε:
$\alpha_2 = 0 \Leftrightarrow \alpha_1 + (2 - 1)\omega = 0 \Leftrightarrow \alpha_1 + \omega = 0 \Leftrightarrow \alpha_1 = -\omega$ (1)
και
\begin{align*}
\alpha_4 &= 4 \Leftrightarrow \alpha_1 + (4 - 1)\omega = 4 \Leftrightarrow \alpha_1 + 3\omega = 4 \stackrel{(1)}{\Leftrightarrow} \\
\Leftrightarrow -\omega + 3\omega &= 4 \Leftrightarrow 2\omega = 4 \Leftrightarrow \omega = 2
\end{align*}

Αντικαθιστούμε στην σχέση (1) και βρίσκουμε:
$\alpha_1 = -2$.

\item Ο $\nu$-οστός όρος της αριθμητικής προόδου είναι:
\begin{align*}
\alpha_\nu = \alpha_1 + (\nu - 1)\omega \ &\Leftrightarrow \\
\Leftrightarrow \alpha_\nu = -2 + (\nu - 1) \cdot 2 \ &\Leftrightarrow \\
\Leftrightarrow \alpha_\nu = -2 + 2\nu - 2 \ &\Leftrightarrow \\
\Leftrightarrow \alpha_\nu &= 2\nu - 4.
\end{align*}

Ισχύει επίσης ότι:
\[\alpha_\nu = 98 \Leftrightarrow 2\nu - 4 = 98 \Leftrightarrow 2\nu = 102 \Leftrightarrow \nu = 51.\]

Άρα ο $51^{\text{ος}}$ όρος της προόδου είναι ίσος με $98$.
\end{alist}

